\documentclass[10pt,a4paper]{amsart}
\usepackage[margin=19mm]{geometry}
\usepackage{amsmath,amssymb,mathtools}
\usepackage[scr=rsfs,cal=euler]{mathalfa}
\usepackage{xfrac}
\usepackage[unicode,hidelinks,bookmarks=false]{hyperref}
\newcommand{\ie}{i.e.\ }
\newcommand{\cf}{cf.\ }
\newcommand{\resp}{respectively\ }
\newcommand{\Subsection}{\subsection}
\providecommand{\nobreakdash}{\nobreak\hbox{-}}
\setlength{\emergencystretch}{2em}
\multlinegap0pt
\usepackage[linesnumbered,ruled,vlined]{algorithm2e}
\usepackage{csquotes}
\newtheorem{theorem}{Theorem}[section]
\newtheorem{example}[theorem]{Example}
\newtheorem{proposition}[theorem]{Proposition}
\newtheorem{corollary}[theorem]{Corollary}
\newtheorem{lemma}[theorem]{Lemma}
\newtheorem{definition}[theorem]{Definition}
\newtheorem{assumption}[theorem]{Assumption}
\newtheorem{observation}[theorem]{Observation}
\title[Revisiting transportation problems under Monge costs]{Revisiting transportation problems under Monge costs with applications to the Discrete Ordered Median Problem: compact optimal dual solutions from the northwest-corner rule (Theorem 2.5)}
\author[Stefan Nickel, Justo Puerto, Simon Ramoser, and Alberto Torrejon]{Stefan Nickel, Justo Puerto, Simon Ramoser, and Alberto Torrejon}
\date{Source: arXiv:2602.15151v3 (CC BY 4.0)}
\begin{document}
\maketitle
\begin{abstract}
We investigate the classical Hitchcock-Koopmans Transportation Problem under the assumption that the cost matrix satisfies the Monge property. In this setting, we derive compact formulas for optimal dual solutions based on the northwest-corner rule. As an application illustrating how these formulas yield structural insight while enhancing computational performance, we consider a broad class of facility location problems. In particular, the expressions are used within a Benders decomposition framework to derive novel mixed-integer linear programming formulations for the Discrete Ordered Median Problem with non-increasing weights. Numerical experiments validate that the resulting formulations achieve state-of-the-art performance and exhibit strong robustness across a wide range of instances.
\end{abstract}

\vspace{1em}
\noindent\textit{Source and licence.} Revisiting transportation problems under Monge costs with applications to the Discrete Ordered Median Problem. Version arXiv:2602.15151v3 (CC BY 4.0), distributed under the Creative Commons Attribution 4.0 International licence, http://creativecommons.org/licenses/by/4.0/, as recorded in the arXiv OAI licence metadata for this exact version and shown on the abstract page https://arxiv.org/abs/2602.15151v3. This item is an editorial re-presentation of the paper's Theorem 2.5 and of the paper's own proof of it. A full rights notice is delivered at licenses/arxiv-2602.15151-CC-BY-4.0.txt.
\noindent\textit{Editorial scope.} Counted unit: the paper's compact closed-form expressions for optimal dual solutions of the balanced transportation problem under Monge costs, printed as Theorem 2.5 (source0.tex lines 284--289), delivered together with the paper's complete original proof of it (source0.tex lines 290--294) as the single primary proof body. No mathematics is written, repaired or re-derived: statement and proof are the paper's own text line for line, in the paper's own notation ($p$, $q$, $s$, $d$, $c_{ij}$, $[p]$, $[q]$, $a^t$, $j_i$, $u$, $v$, $u_i$, $v_j$). Required context retained uncounted: the abstract (source0.tex lines 72--74, reproduced in the wrapper front matter); the heading of the paper's Section 2 with the problem setup paragraph (lines 112 and 115--116); the northwest-corner algorithm as the paper prints it, with the observation on the cells it traverses and that observation's own proof (lines 133--173); the two dual-recursion algorithms (lines 217--232 and 236--251); and the paragraph defining the indices $j_i$ (line 282). Printed numbers are the paper's own: the wrapper restores the paper's section-relative theorem counter and its global equation counter, so the retained section prints as Section 2, the retained observation as Observation 2.1 and the counted statement as Theorem 2.5, and its two displayed formulas print as the paper's (4) and (5). Omitted from this package and not redistributed: the paper's preamble and front matter as such (lines 1--71, of which only the title, the author list and the abstract are reproduced in the wrapper); the introduction and literature review (lines 80--111); the parts of Subsections 2.1--2.4 outside the retained ranges; the whole DOMP application of Section 3 and the computational study and acknowledgements of Section 4; and the paper's own reference list with its bib data file. No citation command occurs in any retained span, so the delivered document carries no bibliography, and no reference command of the delivered text was rewritten. Scope deviation, disclosed: the paper's application theorems (printed as Theorems 3.3 and 3.4) are stated inside optidef model environments whose package is not present in the pinned compile runtime, so they are outside this package; the counted unit is instead the paper's main theoretical result of Section 2, the novel compact dual expressions that the paper's introduction names as its contribution. Printed slips kept literal, among them the paper's own wording, its double-dollar displays, and its reuse of the symbol $i_j$ both for the index defined in the counted proof and in the surrounding text. Layout and class: the paper's own class is the standard article class with the letterpaper option, which is not a publisher class; the delivered wrapper is the lane's pinned amsart editorial wrapper at 10pt on A4, to which this package adds the paper's own nine theorem-like declarations with their shared section-relative counter and its two algorithm support packages. No publisher class file, no publisher style file and no upstream script is redistributed. Assets: the paper contains figure environments and graphics-inclusion commands elsewhere in the paper, but no retained span of this package contains a graphics-inclusion command, no figure environment is retained, and no image file is copied into or redistributed with this package, so all mathematics of the delivered unit is native text with no image dependency. A full rights notice is delivered at licenses/arxiv-2602.15151-CC-BY-4.0.txt.
\setcounter{section}{1}
\section{The Transportation Problem under Monge costs}\label{sec:transportation}

\subsection{Problem definition and the northwest-corner solution method}\label{sec:TP_definition}
Let $p \in \mathbb{Z}_{\geq 1}$ be the number of supply nodes with supplies $s \in \mathbb{R}_{\geq 0}^p$, and $q \in \mathbb{Z}_{\geq 1}$ be the number of demand nodes with demands $d \in \mathbb{R}_{\geq 0}^q$. Moreover, let $c \in \mathbb{R}_{\geq 0}^{p \times q}$ be the transportation cost matrix, where entry $c_{ij}$ prescribes the costs of carrying one unit of the transported commodity from supply node $i$ to demand node $j$. The goal is to find the non-negative quantities that should be transferred between any pair of supply and demand nodes such that all supplies and demands are exactly satisfied and the total transportation cost is minimal. It is easy to see that this is possible if and only if supplies and demands are \enquote{balanced}, so we assume $\sum_{i = 1}^p s_i = \sum_{j = 1}^q d_j$.

\begin{algorithm}[ht]
\caption{Northwest-corner rule}\label{NW_corner_rule}
\DontPrintSemicolon
$x \gets \{\}$\;
$\eta \gets \{\}$\;
$\delta_i^1 \gets s_i \quad \forall \, i \in [p]$ \tcp*[r]{remaining supplies}
$\delta_j^2 \gets d_j \quad \forall \, j \in [q]$ \tcp*[r]{remaining demands}
$a \gets (1,1)$ \tcp*[r]{start in the northwest-corner}
$t \gets 0$\;
\While{$a_2 \leq q$}{
    $t \gets t+1$\;
    $a^t \gets a$\;
    $x_{a^t} \gets \min\{\delta_{a^t_1}^1, \delta_{a^t_2}^2\}$; append $(a^t, x_{a^t})$ to $x$\;
    $\delta_{a^t_1}^1 \gets \delta_{a^t_1}^1 - x_{a^t}$\;
    $\delta_{a^t_2}^2 \gets \delta_{a^t_2}^2 - x_{a^t}$\;
    \If{$\delta_{a^t_1}^1 = 0$ and $a^t_1 < p$}{
        $\eta_t \gets 1$; append $\eta_t$ to $\eta$\;
        $a^t_1 \gets a^t_1+1$ \tcp*[r]{move down}
    }
    \Else{
        $\eta_t \gets 2$ ; append $\eta_t$ to $\eta$\;
        $a^t_2 \gets a^t_2+1$ \tcp*[r]{move right}
    }
    $a \gets a^t$
}
$T \gets t$\;
\Return{$\{(a^1, x_{a^1}), \dots, (a^T, x_{a^T})\}, \{\eta_1, \dots, \eta_{T-1}\}$}
\end{algorithm}

Algorithm~\ref{NW_corner_rule} moves along a staircase-shaped path from cell $(1,1)$ to cell $(p,q)$: between cell $a^t$ and $a^{t+1}$, it either advances one row (\enquote{down}) or one column (\enquote{right}) on the $p \times q$ grid. This movement is well understood and can be characterized as follows.
\begin{observation}\label{NW_corner_characterization}
    We have $(i, j) \in \{ a^1, \dots, a^T \} =: \mathcal{A}$ if and only if one of the following holds:
    \begin{enumerate}
        \item[(i)] $(i, j) = (1, 1)$
        \item[(ii)] $(i-1, j) \in \mathcal{A}$ and $\sum_{k=1}^{i-1} s_k \leq \sum_{k=1}^{j} d_k$ and $i-1 < p$
        \item[(iii)] $(i, j-1) \in \mathcal{A}$ and $\bigl(\sum_{k=1}^{i} s_k > \sum_{k=1}^{j-1} d_k$ or $i = p \bigr)$
    \end{enumerate}
    \begin{proof}
        The algorithm is initialized with $a^1 = (1, 1)$. From there, it proceeds by moving one row downwards or one column to the right. \textit{(ii)} characterizes a downward movement from $(i-1, j)$ to $(i, j)$. This happens exactly when $i - 1 < p$ and the total supply up to $i-1$ does not exceed the total demand up to $j$, that is, when $\sum_{k=1}^{i-1} s_k \leq \sum_{k=1}^{j} d_k$. Similarly, \textit{(iii)} describes a move to the right from $(i, j-1)$ to $(i, j)$. Such a move occurs exactly when the total supply up to $i$ exceeds the total demand up to $j$, that is, $\sum_{k=1}^{i} s_k > \sum_{k=1}^{j-1} d_k$, or when $i = p$. Note that in the latter case, it holds $\sum_{k=1}^{i} s_k = \sum_{k=1}^{j-1} d_k$. Overall, we see that \textit{(i), (ii), (iii)} inductively describe the cells $a^1, \dots, a^T$ traversed by Algorithm~\ref{NW_corner_rule}.
    \end{proof}
\end{observation}

\begin{algorithm}[ht]
\caption{Dual backward recursion}\label{dual_algo_backward}
\DontPrintSemicolon
Declare $u \in \mathbb{R}^p$ and $v \in \mathbb{R}^q$\;
$v_q \gets 0$\;
$\eta_T \gets 1$ \tcp*[r]{$u_p$ will be set in the first iteration}
\For{$t$ from $T$ down to $1$}{
    \If{$\eta_t = 1$}{
        $u_{a_1^t} \gets c_{a^t} - v_{a_2^{t}}$\;
    }
    \ElseIf{$\eta_t = 2$}{
        $v_{a_2^t} \gets c_{a^t} - u_{a_1^{t}}$\;
    }
}
\Return{u, v}
\end{algorithm}

\begin{algorithm}[ht]
\caption{Dual forward recursion}\label{dual_algo_forward}
\DontPrintSemicolon
Declare $u \in \mathbb{R}^p$ and $v \in \mathbb{R}^q$\;
$u_1 \gets 0$\;
$\eta_0 \gets 2$ \tcp*[r]{$v_1$ will be set in the first iteration}
\For{$t$ from $1$ to $T$}{
    \If{$\eta_{t-1} = 1$}{
        $u_{a_1^t} \gets c_{a^t} - v_{a_2^{t}}$\;
    }
    \ElseIf{$\eta_{t-1} = 2$}{
        $v_{a_2^t} \gets c_{a^t} - u_{a_1^{t}}$\;
    }
}
\Return{u, v}
\end{algorithm}

For $i \in \{2, \dots, p\}$, define $j_i := \min \{j \in [q] \, | \, (i,j) \in \{a^1, \dots, a^T\}\}$. As the northwest-corner rule moves in a staircase path from cell $(1,1)$ to cell $(p,q)$ while advancing either one row or one column per iteration, the existence of $j_i$ is guaranteed for all $i \in \{2, \dots, p\}$. Observe that Algorithm~\ref{dual_algo_forward} sets $u_i$ in iteration $t$ where $a^t = (i, j_i)$. Using this notation, the next theorem provides compact formulas for $u$ and $v$ computed by Algorithm~\ref{dual_algo_forward}. This is the first variant of our main result in Section~\ref{sec:transportation}.

\setcounter{theorem}{4}
\setcounter{equation}{3}
\begin{theorem}\label{theorem_u_v_formulas}
    Algorithm~\ref{dual_algo_forward} returns the following solution: \footnote{We assume the convention that an empty sum has value equal to $0$.}
    \begin{align}
        &u_i = \sum_{k=2}^{i} \bigl( c_{k j_k} - c_{k-1, j_k} \bigr) &i \in [p], \label{u_formula} \\
        &v_j = c_{1j} + \sum_{\substack{k \in [p], \, k \geq 2 \\ j_k \leq j}} \bigl( c_{kj} - c_{k-1, j} - c_{k j_k} + c_{k-1, j_k} \bigr) &j \in [q]. \label{v_formula}
    \end{align}

    \begin{proof}
        First, we show \eqref{u_formula} by induction. For $i = 1$, the sum is empty, so $u_1 = 0$ which is coherent with how Algorithm~\ref{dual_algo_forward} sets $u_1$. For the inductive step suppose the formula is correct for some $i \in [p-1]$, that is, $u_i = \sum_{k=2}^{i} \bigl( c_{k j_k} - c_{k-1, j_k} \bigr)$. Recall that, at iteration $t$ with $a^t = (i, j_i)$, Algorithm~\ref{dual_algo_forward} sets $u_{i} \gets c_{i j_i} - v_{j_i}$. It proceeds by setting $v_{j'}$ for all $j' \in \{j_i+1, \dots, j_{i+1}\}$. (If this set is empty, then $j_{i+1} = j_i$ and the procedure continues directly with $u_{i+1}$.) Then, it assigns $u_{i+1} \gets c_{i+1, j_{i+1}} - v_{j_{i+1}}$. But $v_{j_{i+1}} = c_{i j_{i+1}} - u_i$, so $u_{i+1} = u_i + c_{i+1, j_{i+1}} - c_{i j_{i+1}} = \sum_{k = 2}^{i+1} \bigl( c_{k j_k} - c_{k-1, j_k} \bigr)$. We used the induction hypothesis in the last equation. This completes the inductive proof.

        To show \eqref{v_formula}, we make use of the previously shown formula \eqref{u_formula}. In particular, let $j \in [q]$ and $(i_j,j) \in \{a^1, \dots, a^T\}$. Then, Algorithm~\ref{dual_algo_forward} assigns $v_j \gets c_{i_j j} - u_{i_j}$. So, using \eqref{u_formula}, $v_j = c_{i_j j} - \sum_{k=2}^{i_j} \bigl( c_{k j_k} - c_{k-1, j_k} \bigr)$. For $i_j = 1$, it holds that $u_{i_j} = 0$, so the algorithm sets $v_j \gets c_{1j}$ which corresponds to the value given by the formula. For $i_j \geq 2$, observe that $i_j = \max \{k \in [p] \, | \, k \geq 2: j_k \leq j\}$. Therefore, we may rewrite $v_j = c_{i_j j} - \sum_{\substack{k \in [p], \, k \geq 2 \\ j_k \leq j}} \bigl( c_{k j_k} - c_{k-1, j_k} \bigr)$. Moreover, it is easy to see by a telescoping sum argument that $c_{i_j j} = c_{1j} + \sum_{\substack{k \in [p], \, k \geq 2 \\ j_k \leq j}} \bigl( c_{k j} - c_{k-1, j} \bigr)$. Putting everything together, we obtain $v_j = c_{1j} + \sum_{\substack{k \in [p], \, k \geq 2 \\ j_k \leq j}} \bigl( c_{kj} - c_{k-1, j} - c_{k j_k} + c_{k-1, j_k} \bigr)$.
    \end{proof}

\end{theorem}
\end{document}
